\section{Software Interface - Linear Problems}

\subsection{Conventions}

We have adopted the Numerical Recipes in C~\cite{numrecc:88} convention
of using unit offset arrays
where convenient
(almost everywhere).
If an array is declared
{\tt double vect[100];}
then
{\tt (\&vect[0]-1)} or {\tt (vect-1)}
should be passed to
the subroutines below.


\subsection{Subroutine {\tt vspline}}

The syntax for {\tt vspline} is:

\begin{verbatim}

int vspline(ys, alpha, tn, yn, inv_sigma, iid, M, N)
float *alpha;
double *ys, *tn, *yn, *inv_sigma;
int iid, M, N;

\end{verbatim}

%I repeat that the arrays are all unit offset,
%i.e., {\tt alpha[1]} is the first element of {\tt alpha},
%and {\tt alpha[M]} is the last.

{\tt M} is the dimension (length) the measurement vectors $\y_n$.
{\tt N} is the number of samples.
{\tt alpha} is the $M$ smoothing parameters.
{\tt tn} is the $N$ sample points.  If the samples are uniform,
then set {\tt tn} to {\tt (double *) NULL}.
The measurements are stored in 
{\tt yn}
in the order
$\{ y_{1,1},\ldots,y_{1,M}, \ldots, y_{N,1},\ldots,y_{N,M} \}$.
If {\tt iid} is 1, the covariances are assumed to be identical,
and {\tt inv\_sigma} is the $M^2$ elements of $\mcov^{-1}$,
stored by columns or rows
(it is symmetric).
However, if {\tt iid} is 0,
the covariances can vary with $n$,
and {\tt inv\_sigma} is a $N \cdot M^2$ array
of the concatenation of the elements of
$\bSigma_1^{-1}, \ldots, \bSigma_N^{-1}$
in the natural order.

On return,
{\tt ys} is the $N\cdot M$ array of
$\hat{\y}_1,\ldots,\hat{\y}_N$,
stored in the same order as {\tt yn}.
{\tt vspline} returns 1 if all goes well,
0 if there is an error.
The possible errors are: running out of memory,
{\tt tn} not strictly increasing,
or singular covariance matrices.

\subsection{Subroutine {\tt bestcv}}

The syntax for {\tt bestcv} is identical to that of {\tt vspline}.
The difference is that
{\tt alpha} is used as an initial estimate for
$\balpha_{\CV}$.
On return,
{\tt alpha} is set to $\balpha_{\CV}$,
and {\tt ys}
is the smoothed estimate at that value of $\balpha$.

\subsection{Subroutine {\tt cubic\_interp}}

Each component function of
$\bghata(t)$
is a scalar cubic spline.
To compute values of
$\hat{g}_{\balpha,m}(t)$
for values of $t$ other than the knots,
call 
{\tt cubic\_interp}.
%Currently,
%{\tt cubic\_interp}
%is implemented only for scalar functions.
%The syntax is:

\begin{verbatim}

int cubic_interp(yk, tn, ys, tk, N, K)
double *yk, *tn, *ys, *tk;
int N, K;

\end{verbatim}

{\tt tn} is the $N$ sample points,
and {\tt ys} is the $N$ smoothed estimates
for one of the component functions,
e.g.
$\{\hat{y}_{n,m}\}_{n=1}^N$. \
{\tt tk} is $K$ new points at which
$\hat{g}_{\balpha,m}(t)$
is to be evaluated.
On return,
{\tt yk} is the $K$ interpolated samples.
{\tt cubic\_interp} returns 1 if successful, 0 otherwise.
{\tt cubic\_interp} is not very efficient,
and may not be exact for values of $t$
outside of $[t_1,t_N]$.
